\section*{Bab \ref{ch:b2:genfun} --- Fungsi Pembangkit Peluang}

\begin{solution}{exo:b2:genfun:1}
Dadu setimbang: $G(t) = \frac16(t + t^2 + \dots + t^6) = \frac t6(1 + t
+ \dots + t^5)$. Andaikan jumlah dua dadu setimbang seragam pada
$\{2, \dots, 12\}$, maka
\[
G(t)^2 = \frac{t^2}{36}\,h(t)^2
= \frac{t^2}{11}\sum_{k=0}^{10}t^k ,
\qquad h(t) = 1 + t + \dots + t^5 .
\]
Padahal $h$ merupakan suku banyak real berderajat ganjil $5$, jadi ia punya akar
real (teorema nilai antara; secara nyata $h(-1) = 0$), sehingga
$h^2$ punya akar real. Namun $\sum_{k=0}^{10}t^k$ tak punya: ia
positif untuk $t \geq 0$, dan untuk $t < 0$ ia sama dengan
$\frac{t^{11} - 1}{t - 1}$, yakni hasil bagi dua bilangan negatif.
Kontradiksi --- jadi jumlah dua dadu setimbang tak pernah seragam (sebagaimana
dibenarkan \omterm{def:b2:randomvar:law}{distribusi} segitiga jumlah dadu yang sudah kita kenal).
\end{solution}

\begin{solution}{exo:b2:genfun:2}
\emph{Binomial:} $G(t) = (1 - p + pt)^n$, $G'(t) = np(1 - p +
pt)^{n-1}$, $G''(t) = n(n-1)p^2(1 - p + pt)^{n-2}$, jadi
\[
\E(X) = G'(1) = np,
\qquad
V(X) = G''(1) + G'(1) - G'(1)^2
= n(n-1)p^2 + np - n^2p^2 = np(1-p).
\]
\emph{Geometrik} ($q = 1 - p$): $G(t) = \frac{pt}{1 - qt}$, jadi
$G'(t) = \frac{p}{(1 - qt)^2}$ dan $G''(t) = \frac{2pq}{(1 -
qt)^3}$; di $t = 1$ (dengan memakai $1 - q = p$):
\[
\E(X) = \frac{p}{p^2} = \frac1p,
\qquad
V(X) = \frac{2q}{p^2} + \frac1p - \frac{1}{p^2}
= \frac{2q + p - 1}{p^2}
= \frac{q}{p^2} ,
\]
yang cocok dengan \cref{exo:b2:randomvar:1} dengan kerja yang lebih sedikit.
\end{solution}

\begin{solution}{exo:b2:genfun:3}
Tidak, bahkan dengan kecurangan yang berbeda. Andaikan $X, Y$ \omterm{def:b2:randomvar:law}{hukum} pada
$\{1, \dots, 6\}$ yang jumlahnya seragam. Maka $G_X(t) = t\,a(t)$ dan
$G_Y(t) = t\,b(t)$ dengan $a, b$ suku banyak real berderajat \emph{paling
besar} $5$ --- dan derajatnya wajib berjumlah $10$ (jumlahnya mencapai
$12$ dengan peluang positif), jadi $\deg a = \deg b = 5$, keduanya
ganjil. Seperti pada \cref{exo:b2:genfun:1},
\[
a(t)\,b(t) = \frac{1}{11}\sum_{k=0}^{10}t^k
\]
akan memaksa adanya akar real di ruas kirinya (setiap suku banyak real
berderajat ganjil punya satu) dan tak ada di ruas kanannya. Jadi tak ada kecurangan atas
dua dadu yang \omterm{def:b2:proba:independence}{saling bebas} --- sama atau tidak --- yang menghasilkan jumlah seragam.
\end{solution}

\begin{solution}{exo:b2:genfun:4}
Menurut \cref{thm:b2:genfun:compound} dengan $G_N(s) = e^{\lambda(s-1)}$
dan $G_X(t) = 1 - p + pt$:
\[
G_S(t) = e^{\lambda(1 - p + pt - 1)} = e^{\lambda p(t - 1)} :
\]
$S \sim \mathcal{P}(\lambda p)$. Adapun cacah yang terbuang $D = N - S$
mencacah benda yang sama yang disimpan dengan peluang $1 - p$, jadi menurut
perhitungan yang sama $D \sim \mathcal{P}(\lambda(1 - p))$.
Kesalingbebasannya, secara langsung: untuk $j, k \in \N$,
\begin{align*}
\P(S = j,\ D = k)
&= \P(N = j + k)\,\binom{j+k}{j}p^jq^k
= e^{-\lambda}\frac{\lambda^{j+k}}{(j+k)!}\,
\frac{(j+k)!}{j!\,k!}\,p^jq^k\\
&= \Bigl(e^{-\lambda p}\frac{(\lambda p)^j}{j!}\Bigr)
\Bigl(e^{-\lambda q}\frac{(\lambda q)^k}{k!}\Bigr)
\end{align*}
dengan $q = 1 - p$: \omterm{def:b2:randomvar:law}{hukum} gabungannya terfaktorkan menjadi
$\mathcal{P}(\lambda p) \otimes \mathcal{P}(\lambda q)$. Sebuah arus Poisson
yang terbelah secara acak menghasilkan arus Poisson yang \emph{\omterm{def:b2:proba:independence}{saling bebas}} ---
mukjizat kecil yang terus-menerus dipakai dalam teori antrean.
\end{solution}

\begin{solution}{exo:b2:genfun:5}
Waktu tunggu antara gambar yang berurutan merupakan peubah geometrik
$\mathcal{G}(p)$ yang \omterm{def:b2:proba:independence}{saling bebas} (sifat tanpa ingatannya: sesudah setiap
gambar permainannya mulai lagi), jadi $T_r = W_1 + \dots + W_r$ dan
kemultiplikatifannya (\cref{thm:b2:genfun:product}) memberi
\[
G_{T_r}(t) = \Bigl(\frac{pt}{1 - qt}\Bigr)^{r},
\qquad
\E(T_r) = r\,\E(W_1) = \frac rp,
\qquad
V(T_r) = r\,V(W_1) = \frac{rq}{p^2}
\]
($q = 1 - p$; \omterm{def:b2:randomvar:variance}{ragamnya} menjumlah menurut kesalingbebasannya). Penguraiannya: menurut
deret binomial yang diperumum
(\cref{ch:b2:powerseries}), $(1 - qt)^{-r} = \sum_{m\geq0}
\binom{m + r - 1}{r - 1}q^mt^m$, jadi koefisien $t^n$ pada
$p^rt^r(1 - qt)^{-r}$ adalah (dengan $m = n - r$)
\[
\P(T_r = n) = \binom{n-1}{r-1}p^r(1-p)^{n-r},
\qquad n \geq r ,
\]
yakni \omterm{def:b2:randomvar:law}{hukum} \emph{binomial negatif} --- secara kombinatorik: gambar ke-$r$
jatuh pada lemparan $n$ bila dan hanya bila $r - 1$ gambar sebelumnya memilih
tempatnya di antara $n - 1$ lemparan pertamanya.
\end{solution}

\begin{solution}{exo:b2:genfun:6}
$m = 1\cdot\frac38 + 2\cdot\frac38 + 3\cdot\frac18 = \frac{3 + 6 +
3}{8} = \frac32 > 1$: superkritis. \omterm{ex:b2:powerseries:fibonacci}{Fungsi pembangkitnya} adalah
\[
G(t) = \frac{1 + 3t + 3t^2 + t^3}{8} = \frac{(1 + t)^3}{8} ,
\]
jadi titik tetapnya memecahkan $(1 + t)^3 = 8t$, yakni $t^3 + 3t^2 - 5t + 1
= 0$. Dengan memfaktorkan keluar akar terjaminnya $t = 1$:
\[
t^3 + 3t^2 - 5t + 1 = (t - 1)\bigl(t^2 + 4t - 1\bigr),
\]
dan $t^2 + 4t - 1 = 0$ memberi $t = -2 \pm \sqrt5$. Akar pada
$\intco{0}{1}$ adalah $\sqrt5 - 2 \approx 0.236$: menurut
\cref{thm:b2:genfun:extinction},
\[
q = \sqrt 5 - 2 .
\]
(Pemeriksaan yang menyenangkan: \omterm{def:b2:randomvar:law}{hukum} keturunannya adalah \omterm{def:b2:randomvar:law}{hukum} $3$ koin
setimbang yang \omterm{def:b2:proba:independence}{saling bebas}, $Z_1 \sim \mathcal{B}(3, \frac12)$.)
\end{solution}

\begin{solution}{exo:b2:genfun:7}
\emph{\omterm{def:b2:randomvar:expectation}{Nilai harapannya}.} Mula-mula $\E(Z_n) = m^n$: menurut
\cref{thm:b2:genfun:compound}, $\E(Z_{n+1}) = \E(Z_n)\,m$, dan
$\E(Z_0) = 1$. Keluarga $\bigl(Z_n(\omega)\P(\{\omega\})
\bigr)_{n, \omega}$ taknegatif, jadi Fubini untuk keluarga berlaku
tanpa syarat:
\[
\E(Y) = \sum_{n=0}^{\infty}\E(Z_n)
= \sum_{n=0}^\infty m^n = \frac{1}{1 - m} < \infty
\]
(khususnya $Y$ berhingga hampir pasti: sesuai dengan
kepunahan yang pasti pada kasus subkritisnya).

\emph{Persamaan fungsionalnya.} Uraikan populasinya menurut
anak leluhurnya: bila leluhurnya punya $Z_1 = k$ anak, maka
total keturunannya adalah $Y = 1 + Y_1 + \dots + Y_k$, dengan $Y_i$ menyatakan
total keturunan garis anak ke-$i$ --- dan $Y_i$ merupakan
salinan $Y$ yang \omterm{def:b2:proba:independence}{saling bebas}, bebas pula dari $Z_1$ (garis yang berbeda
memakai \omterm{def:b2:proba:space}{kejadian} perkembangbiakan yang saling lepas dan \omterm{def:b2:proba:independence}{saling bebas}). Dengan menyaratkan pada
$Z_1$ seperti pada \cref{thm:b2:genfun:compound}:
\[
H(t) = \E\bigl(t^Y\bigr)
= t\sum_{k=0}^\infty \P(Z_1 = k)\,H(t)^k
= t\,G\bigl(H(t)\bigr),
\]
dengan faktor $t$ memperhitungkan leluhurnya sendiri. (Untuk
\omterm{def:b2:randomvar:law}{hukum} $p_0 = 1 - p$, $p_2 = p$ pada percabangan biner, persamaan kuadrat
dalam $H$ ini dapat dipecahkan secara eksplisit lalu diuraikan --- dan
\omterm{ex:b2:powerseries:catalan}{bilangan Catalan} pada \cref{ch:b2:powerseries} mencacah pohon
keluarganya.)
\end{solution}

\begin{solution}{exo:b2:genfun:8}
Misalkan $R > 1$ jari-jarinya lalu tetapkan $t \in \intoo{1}{R}$. Maka
$\E(t^X) = G(t) < \infty$, dan ketaksamaan Markov
(\cref{thm:b2:randomvar:markov}) yang diterapkan pada peubah taknegatif
$t^X$ pada aras $t^n$:
\[
\P(X \geq n) = \P\bigl(t^X \geq t^n\bigr)
\leq \frac{G(t)}{t^n} = C\rho^n,
\qquad C = G(t),\quad \rho = \frac1t \in \intoo{0}{1}.
\]
\emph{Konversnya:} bila $\P(X \geq n) \leq C\rho^n$, maka $p_n \leq
\P(X \geq n) \leq C\rho^n$, jadi untuk $\abs t < \frac1\rho$ deret
$\sum p_n\abs t^n$ terdominasi oleh deret geometrik yang konvergen
$C\sum(\rho\abs t)^n$: jari-jarinya sedikitnya
$\frac1\rho > 1$. Jari-jari \omterm{ex:b2:powerseries:fibonacci}{fungsi pembangkitnya} dan peluruhan geometrik
ekornya merupakan dua muka satu sifat yang sama.
\end{solution}

\begin{solution}{exo:b2:genfun:9}
Tuliskan $p_k^{(n)} = \P(X_n = k)$, $p_k = \P(X = k)$.

\emph{Kasus $k = 0$.} Untuk $t \in \intoo{0}{1}$ dan sebarang \omterm{def:b2:randomvar:law}{hukum}
$(q_j)$ dengan $\sum_j q_j \leq 1$:
\[
\Bigl|\,q_0 - \sum_j q_jt^j\Bigr|
= \sum_{j \geq 1} q_j t^j
\leq \sum_{j\geq1}t^j = \frac{t}{1 - t} .
\]
Jadi
\[
\abs{p_0^{(n)} - p_0}
\leq \frac{2t}{1 - t}
+ \abs{G_{X_n}(t) - G_X(t)} .
\]
Diberikan $\varepsilon > 0$, pilihlah $t$ dengan $\frac{2t}{1-t} <
\frac\varepsilon2$, lalu $n_0$ sedemikian sehingga suku terakhirnya $<
\frac\varepsilon2$ untuk $n \geq n_0$: jadi $p_0^{(n)} \to p_0$.

\emph{Langkah induksinya.} Andaikan $p_j^{(n)} \to p_j$ untuk $j < k$.
Tinjaulah fungsi yang \emph{tergeser}
\[
g_n(t) = \frac{G_{X_n}(t) - p^{(n)}_0}{t}
= \sum_{j\geq0} p^{(n)}_{j+1}t^j,
\qquad
g(t) = \frac{G_X(t) - p_0}{t} ,
\]
yakni \omterm{ex:b2:powerseries:fibonacci}{fungsi pembangkit} barisan subpeluang
$(p^{(n)}_{j+1})_j$ (bermassa total $\leq 1$, dan hanya itulah yang dipakai
hujah $k = 0$). Untuk $t \in \intoo{0}{1}$ yang tetap, $g_n(t)
\to g(t)$ menurut hipotesisnya dan kasus $k = 0$. Menerapkan hujah $k =
0$ pada $g_n$ memberi $p_1^{(n)} \to p_1$; dengan mengiterasi
geserannya $k$ kali diperoleh $p_k^{(n)} \to p_k$ untuk setiap $k$. (Inilah
contoh diskret dan dasar teorema kekontinuan Lévy,
yang bentuk umumnya --- bagi fungsi karakteristik --- merupakan tonggak Tahun
ke-3.)
\end{solution}

\begin{solution}{exo:b2:genfun:10}
Secara \omterm{def:b2:funcseq:def}{titik demi titik}, $\frac{1 + (-1)^X}{2}$ bernilai $1$ ketika $X$ genap
dan $0$ ketika ganjil, jadi dengan mengambil \omterm{def:b2:randomvar:expectation}{nilai harapannya} (pemindahan),
\[
\P(X \text{ genap}) = \frac{1 + \E\bigl((-1)^X\bigr)}2 =
\frac{1 + G_X(-1)}2 .
\]
Poisson: $\frac{1 + \eu^{-2\lambda}}2 \to \frac12$ ketika
$\lambda$ membesar. Binomial: $\frac{1 + (1 - 2p)^n}2$. Pada kedua
kasusnya $G_X(-1) \to 0$ berarti keparitasan $X$ menjadi koin
yang setimbang: \omterm{def:b2:randomvar:law}{hukumnya} menyebar ke banyak bilangan bulat lalu melupakan
keparitasannya.
\end{solution}

\begin{solution}{exo:b2:genfun:11}
$t + \dots + t^6 = t\,\frac{1 - t^6}{1 - t}$ dan $1 - t^6 =
(1 - t)(1 + t)(1 + t + t^2)(1 - t + t^2)$, yang memberikan pemfaktoran
yang dinyatakan itu. Untuk dadu pertamanya, $(1 + t)(1 + t + t^2) = 1 +
2t + 2t^2 + t^3$, jadi $\frac{t(1+t)(1+t+t^2)}6 = \frac{t +
2t^2 + 2t^3 + t^4}6$: bermuka $\{1, 2, 2, 3, 3, 4\}$. Untuk yang
kedua, dengan menguraikan
\[
(1 + 2t + 2t^2 + t^3)(1 - t + t^2)^2 = 1 + t^2 + t^3 + t^4 +
t^5 + t^7,
\]
jadi $\frac{t(1+t)(1+t+t^2)(1-t+t^2)^2}6 = \frac{t + t^3 + t^4
+ t^5 + t^6 + t^8}6$: bermuka $\{1, 3, 4, 5, 6, 8\}$. Adapun
hasil kali kedua \omterm{ex:b2:powerseries:fibonacci}{fungsi pembangkitnya} menyusun ulang keenam
faktornya menjadi $\bigl(\frac{t(1+t)(1+t+t^2)(1-t+t^2)}6
\bigr)^2$, yakni kuadrat fungsi dadu bakunya: pasangan
Sicherman itu berhukum persis sama dengan yang baku bagi totalnya ---
dan \omterm{ex:b2:powerseries:fibonacci}{fungsi pembangkit} mengelompokkan semua penyusunan ulang semacam itu.
\end{solution}

\begin{solution}{exo:b2:genfun:12}
Misalkan $A$ dan $B$ \omterm{ex:b2:powerseries:fibonacci}{fungsi pembangkit} sisa
lamanya permainan yang bermula dari ``belum ada gambar'' dan ``sudah satu
gambar''. Satu lemparan terpakai, lalu: dari keadaan $0$, angka
mengembalikannya ke keadaan $0$, sedangkan gambar memindahkannya ke keadaan $1$; dari keadaan
$1$, gambar mengakhiri permainannya, sedangkan angka mengembalikannya ke keadaan $0$:
\[
A(t) = t\bigl(q\,A(t) + p\,B(t)\bigr),
\qquad
B(t) = t\bigl(p + q\,A(t)\bigr).
\]
Dengan menyubstitusikannya: $A(1 - qt) = pt\,B = pt(pt + qtA)$, jadi
\[
G_T(t) = A(t) = \frac{p^2t^2}{1 - qt - pq\,t^2} .
\]
Di $t = 1$ penyebutnya adalah $1 - q - pq = p(1 - q) = p^2$:
$G_T(1) = 1$, jadi permainannya berakhir hampir pasti (sebagaimana
ditunjukkan \cref{exo:b2:proba:6} lewat rekursi). Penurunan
logaritmik di $1$: $\E(T) = 2 - \frac{D'(1)}{D(1)}$ dengan
$D(t) = 1 - qt - pqt^2$, $D'(1) = -q - 2pq$:
\[
\E(T) = 2 + \frac{q + 2pq}{p^2} = \frac{2p^2 + q + 2pq}{p^2}
= \frac{1 + p}{p^2},
\]
yang bernilai $6$ untuk $p = \frac12$.
\end{solution}

\begin{solution}{pb:b2:genfun:1}
\textbf{1.} Untuk $t \in \intoo01$, aturan rantai pada $G_n = G
\circ G_{n-1}$ memberi $G_n'(t) =
G'\bigl(G_{n-1}(t)\bigr)G_{n-1}'(t)$. Ketika $t \to 1^-$,
$G_{n-1}(t) \uparrow 1$, dan $G'$ tidak turun dengan
limit kiri $m$ di $1$, jadi faktor pertamanya menuju $m$; menurut
induksi faktor keduanya menuju $m^{n-1}$. Menurut
\cref{thm:b2:genfun:moments}, $\E(Z_n) = G_n'(1^-) = m^n$.

\textbf{2.} Dengan menurunkannya sekali lagi,
\[
G_n'' = G''(G_{n-1})\,(G_{n-1}')^2 +
G'(G_{n-1})\,G_{n-1}'',
\]
lalu dengan membiarkan $t \to 1^-$: $a_n = G''(1)m^{2(n-1)} + m\,
a_{n-1}$ dengan $a_n = G_n''(1)$, $a_1 = G''(1)$. Untuk $m \neq
1$ orang memeriksa menurut induksi bahwa $a_n = G''(1)\,m^{n-1}
\frac{m^n - 1}{m - 1}$ (rekursinya menambahkan $G''(1)m^{2n-2}$
pada $m\cdot G''(1)m^{n-2}\frac{m^{n-1}-1}{m-1}$, dan
$m^{n-1} + \frac{m^{n-1}-1}{m-1} = \frac{m^n - 1}{m-1}$);
untuk $m = 1$, $a_n = a_{n-1} + G''(1) = n\,G''(1)$.

\textbf{3.} $V(Z_n) = a_n + m^n - m^{2n}$ dan $G''(1) =
\sigma^2 + m^2 - m$. Untuk $m \neq 1$, potongan $(m^2 -
m)m^{n-1}\frac{m^n-1}{m-1} = m^n(m^n - 1)$ meniadakan $m^n -
m^{2n}$ secara persis, sehingga tersisa $V(Z_n) =
\sigma^2m^{n-1}\frac{m^n-1}{m-1}$. Untuk $m = 1$: $V(Z_n) =
nG''(1) = n\sigma^2$.

\textbf{4.} $Z_n$ merupakan peubah bulat taknegatif, jadi
$\P(Z_n > 0) = \P(Z_n \geq 1) \leq \E(Z_n) = m^n$ menurut Markov
(\cref{thm:b2:randomvar:markov}). Untuk $m < 1$ ini meluruh
secara geometrik --- dan \omterm{def:b2:series:summable}{terjumlahkan}, jadi Borel--Cantelli bahkan memberi
bahwa hanya berhingga banyak generasinya yang takkosong, dan itu pun
kepunahan lagi.

\textbf{5.} Cauchy--Schwarz: $\E(Z_n)^2 = \E(Z_n\mathbf
1_{Z_n>0})^2 \leq \E(Z_n^2)\,\P(Z_n > 0)$. Dengan pertanyaan 3
dan $m < 1$:
\[
\E(Z_n^2) = V(Z_n) + m^{2n}
\leq \frac{\sigma^2m^{n-1}}{1-m} + m^{2n},
\]
jadi, dengan membagi $m^{2n}$ oleh batas ini lalu menyederhanakannya oleh
$m^n$,
\[
\P(Z_n > 0) \geq \frac{m^n}{\frac{\sigma^2}{m(1-m)} + m^n}
\geq \Bigl(\frac{\sigma^2}{m(1-m)} + 1\Bigr)^{-1}m^n ,
\]
dengan memakai $m^n \leq 1$ pada penyebutnya. Dengan pertanyaan 4:
$\P(Z_n > 0) \asymp m^n$.

\textbf{6.} $G(t) = q\sum_k(pt)^k = \frac{q}{1 - pt}$, dan $m
= G'(1) = \frac{pq}{(1-p)^2} = \frac pq$. Subkritis untuk $p
< \frac12$, kritis untuk $p = \frac12$, dan superkritis untuk $p
> \frac12$.

\textbf{7.} $G(t) = t$ berbunyi $pt^2 - t + q = 0$, dengan akar
$\frac{1 \pm \abs{p - q}}{2p}$, yakni $1$ dan $\frac qp =
\frac1m$. Peluang kepunahannya adalah titik tetap terkecil
pada $\intcc01$ (\cref{thm:b2:genfun:extinction}):
$q_{\mathrm{ext}} = 1$ bila $m \leq 1$, dan $\frac1m$ bila $m >
1$.

\textbf{8.} Untuk $m \neq 1$, dengan $p = \frac m{m+1}$, $q =
\frac1{m+1}$: bila $q_n = \frac{m^n - 1}{m^{n+1} - 1}$, maka
\[
1 - p\,q_n = \frac{(m+1)(m^{n+1} - 1) - m(m^n - 1)}
{(m+1)(m^{n+1} - 1)} = \frac{m^{n+2} - 1}{(m+1)(m^{n+1} -
1)},
\]
jadi $q_{n+1} = \frac{q}{1 - pq_n} = \frac{m^{n+1} -
1}{m^{n+2} - 1}$; dan kasus dasarnya $q_0 = 0$ berlaku. Untuk $m = 1$:
$G(t) = \frac1{2 - t}$ dan $q_{n+1} = \frac1{2 -
\frac{n}{n+1}} = \frac{n+1}{n+2}$, dengan $q_0 = 0$.

\textbf{9.} $1 - q_n = \frac{m^n(m - 1)}{m^{n+1} - 1}$. Untuk
$m < 1$ penyebutnya menuju $-1$: $1 - q_n \sim (1 -
m)\,m^n$. Untuk $m > 1$:
\[
q_{\mathrm{ext}} - q_n = \frac1m - \frac{m^n - 1}{m^{n+1} -
1} = \frac{m - 1}{m\,(m^{n+1} - 1)} \sim \frac{m -
1}{m^{2}}\;m^{-n} .
\]
Adapun $G'(t) = \frac{pq}{(1 - pt)^2}$ yang dievaluasi di $t = \frac
qp$ (tempat $1 - pt = 1 - q = p$) memberi $G'(q_{\mathrm{ext}})
= \frac qp = \frac1m$: nisbah teramatinya $m^{-1}$ persis
merupakan turunannya di titik tetap penariknya.

\textbf{10.} Untuk $p = \frac12$: $G''(t) = \frac{1/4}{(1 -
t/2)^3}$, jadi $G''(1) = 2$ dan $\sigma^2 = G''(1) + m - m^2 =
2$. \omterm{def:b2:multint:exact}{Bentuk tertutupnya} memberi $1 - q_n = \frac1{n+1}$: adapun
peluang kesintasannya meluruh seperti $1/n$ --- terlalu pelan untuk
\omterm{def:b2:series:summable}{terjumlahkan}, tak seperti laju subkritis mana pun.

\textbf{11.} Induksi: $G_1(t) = \frac1{2-t}$ cocok dengan
rumusnya untuk $n = 1$, dan
\[
G(G_n(t)) = \cfrac{1}{2 - \cfrac{n - (n-1)t}{n+1 - nt}}
= \frac{n + 1 - nt}{2(n+1) - 2nt - n + (n-1)t}
= \frac{n+1 - nt}{n + 2 - (n+1)t} .
\]
Lalu
\[
\frac{G_n(t) - q_n}{1 - q_n}
= (n+1)\,\Bigl(\frac{n - (n-1)t}{n+1 - nt} -
\frac{n}{n+1}\Bigr)
= \frac{t}{n + 1 - nt}
= \frac{\frac{t}{n+1}}{1 - \frac{n}{n+1}t} ,
\]
yakni \omterm{ex:b2:powerseries:fibonacci}{fungsi pembangkit} \omterm{def:b2:randomvar:law}{hukum} geometrik $\mathcal G\bigl(\frac1{n+1}
\bigr)$ pada $\N^*$ (\cref{ex:b2:genfun:classical}): bila
kesintasannya diketahui, $\P(Z_n = k \mid Z_n > 0) =
\frac1{n+1}\bigl(\frac n{n+1}\bigr)^{k-1}$, dengan rerata
bersyarat $n + 1$. Adapun rerata tanpa syaratnya $1 = \E(Z_n)$ merupakan
hasil kali peluang kesintasan yang melenyap dan ukuran bersyarat yang
tumbuh secara linear.

\textbf{12.} Bila garisnya punah pada generasi $n$, maka
$Y = Z_0 + \dots + Z_{n-1}$ berhingga; bila ia tak pernah
punah, maka $Y \geq \sum_n 1 = \infty$. Jadi $\{Y < \infty\}$ merupakan
\omterm{def:b2:proba:space}{kejadian} kepunahannya dan $\P(Y < \infty) = q_{\mathrm{ext}}$.
Adapun penurunan $H(t) = tG(H(t))$ pada
\cref{exo:b2:genfun:7} --- leluhurnya menyumbang
faktor $t$, sedangkan anaknya mendirikan salinan $Y$ yang \omterm{def:b2:proba:independence}{saling bebas} dan
tercacah lewat $G$ --- hanya memakai
\cref{thm:b2:genfun:compound}, yang sah pada setiap rezimnya.

\textbf{13.} Dengan $G(s) = \frac{1 + s^2}2$ persamaannya berbunyi
$tH^2 - 2H + t = 0$, jadi $H = \frac{1 - \sqrt{1 - t^2}}{t}$
(yakni akar dengan $H(0) = 0$). Dengan membandingkannya dengan deret Catalan
$C(x) = \frac{1 - \sqrt{1 - 4x}}{2x}$
(\cref{ex:b2:powerseries:catalan}): $H(t) = \frac
t2\,C\bigl(\frac{t^2}4\bigr) =
\sum_{k\geq0}C_k\,\frac{t^{2k+1}}{2^{2k+1}}$, yakni $\P(Y =
2k+1) = C_k2^{-2k-1}$. Pemeriksaannya: $\P(Y = 1) = C_0/2 = \frac12$
(leluhurnya tak punya anak); $\P(Y = 3) = C_1/8 = \frac18$
(dua anak, keduanya tanpa anak: $\frac12\cdot\frac12\cdot
\frac12$).

\textbf{14.} Dengan menurunkan $H = tG(H)$ pada $\intoo01$ lalu
membiarkan $t \to 1^-$ (limit monotonnya seperti pada
\cref{thm:b2:genfun:moments}): $H'(1)\bigl(1 - G'(H(1))\bigr)
= G(H(1))$. Pada kasus subkritisnya $H(1) = 1$ dan $\E(Y) =
H'(1) = \frac1{1 - m}$. Pada kasus kritisnya $G'(1) = 1$
membuat faktor kirinya lenyap sedangkan ruas kanannya bernilai $1$: tak ada
$H'(1)$ berhingga yang mungkin ada, jadi $\E(Y) = \infty$ --- padahal $\P(Y <
\infty) = q = 1$.

\textbf{15.} $C_k = \frac1{k+1}\binom{2k}k \sim
\frac{4^k}{\sqrt\pi\,k^{3/2}}$ menurut
\cref{ex:b2:comparison:centralbinomial}, jadi
\[
\P(Y = 2k+1) = \frac{C_k}{2\cdot4^{k}} \sim
\frac1{2\sqrt\pi\,k^{3/2}} .
\]
Dengan menjumlahkan ekornya (perbandingan dengan $\int_K^\infty
k^{-3/2}\dd k = 2K^{-1/2}$, dari atas dan dari bawah): $\P(Y > 2K)
\asymp K^{-1/2}$, yakni $\P(Y > n) \asymp n^{-1/2}$ --- sebuah
ekor berat dengan rerata tak hingga, yang mengkuantifikasi pertanyaan 14.

\textbf{16.} Kedua objek kritisnya --- waktu pulang jalan
setimbang (soal akhir pekan pada \cref{ch:b2:proba}) dan
total keturunan kritisnya --- berhingga hampir pasti dengan
rerata tak hingga, dengan \omterm{def:b2:randomvar:law}{hukum} lokal bereksponen $-3/2$ dan ekor
bereksponen $-1/2$. Ini bukan kebetulan: menjelajahi sebuah
pohon keluarga anak demi anak akan menghasilkan lintasan $\pm1$ (satu langkah
naik per kelahiran, satu langkah turun per kematian) yang persis merupakan jalan
setimbang, dan $Y$ menjadi waktu lintas pertamanya. Kekritisannya berarti
hanyutan nol: prosesnya selalu berada di ambang kepunahan
sekaligus ledakan, dan fluktuasi berskala $\sqrt{}$
dari keacakan yang berhanyutan nol menghasilkan persis eksponen
ini.

\textbf{17.} $G''(t) = \sum_{n\geq2}n(n-1)p_nt^{n-2}$ bersuku
taknegatif, jadi ia tidak turun pada $\intco01$ dengan
limit berhingga $G''(1) = \sigma^2$ (kekritisannya membuat
$\E Z_1(Z_1 - 1) = \sigma^2$); dan fungsi yang tidak turun dengan
limit yang sama dengan nilai perbatasannya itu \omterm{def:b2:metric:continuity}{kontinu} di $1$.
Taylor dengan sisa integral di titik $1$:
\[
G(t) = 1 + (t - 1) + \int_1^t(t - s)G''(s)\,\dd s
= t + \frac{G''(1)}2(1-t)^2 + o\bigl((1-t)^2\bigr),
\]
sebab $G''(s) = G''(1) + o(1)$ ketika $s \to 1^-$.

\textbf{18.} Dengan menyamakan penyebutnya, $h(t) =
\frac{G(t) - t}{(1 - G(t))(1 - t)}$. Menurut pertanyaan 17
pembilangnya adalah $b(1-t)^2 + o((1-t)^2)$ dan $1 - G(t) = (1 -
t)\bigl(1 - b(1-t) + o(1-t)\bigr)$, jadi $h(t) \to b$.

\textbf{19.} Menurut definisi $h$ di $t = q_j$ dan $G(q_j) =
q_{j+1}$: $\frac1{1 - q_{j+1}} - \frac1{1-q_j} = h(q_j)$;
menjumlahkannya dari $j = 0$ ($q_0 = 0$) memberikan paparannya. Karena
proses kritisnya punah, $q_j \uparrow 1$, jadi $h(q_j) \to
b$ dan rerata Cesàro $\frac1n\sum_{j<n}h(q_j) \to b$:
$\frac1{1-q_n} \sim bn$, yakni
\[
\P(Z_n > 0) \sim \frac1{bn} = \frac{2}{\sigma^2 n} .
\]

\textbf{20.} Kasus geometrik kritisnya: $\sigma^2 = 2$
(pertanyaan 10), jadi Kolmogorov meramalkan $1 - q_n \sim \frac1n$
--- dan nilai persisnya adalah $\frac1{n+1}$.

\textbf{21.} Karena $Z_n\mathbf 1_{Z_n > 0} = Z_n$, maka $\E(Z_n
\mid Z_n > 0) = \frac{\E(Z_n)}{\P(Z_n > 0)} = \frac1{1 -
q_n} \sim \frac{\sigma^2n}2$. Pada kasus geometriknya ini bernilai
$n + 1$, yang cocok persis dengan pertanyaan 11 ($\sigma^2 = 2$). Adapun
gambaran kritisnya: kepunahannya pasti, ukuran reratanya
membeku di $1$, dan garis sintas yang langka itu berukuran tumbuh
secara linear --- masing-masing faktornya mengimbangi yang lain.

\textbf{22.} Untuk keturunan $\mathcal P(\lambda)$, $G(t) =
\eu^{\lambda(t-1)}$ dan peluang kepunahannya adalah
akar terkecil $q = \eu^{\lambda(q-1)}$. Secara numerik:
$\lambda = 1.5$ memberi $q \approx 0.417$ (iterasikan $q \mapsto
\eu^{1.5(q-1)}$: $0, 0.223, 0.312, 0.356, \dots \to 0.4172$);
$\lambda = 2$ memberi $q \approx 0.203$. Jadi satu kasus indeks
memantik wabah besar dengan peluang $58\%$
($\lambda = 1.5$) atau $80\%$ ($\lambda = 2$) --- mungkin, tetapi tak
pasti. Iterasi dari $q_0 = 0$ konvergen ke
akar \emph{terkecilnya} sebab $G$ tidak turun: menurut
induksi $q_n \leq r$ bagi sebarang titik tetap $r$, dan $(q_n)$
naik (ia adalah $\P(Z_n = 0)$), jadi limitnya merupakan titik tetap
di bawah semua yang lain.

\textbf{23.} Ke-$k$ leluhurnya mendirikan pohon keluarga
yang \omterm{def:b2:proba:independence}{saling bebas}, dan kepunahan totalnya adalah irisan $k$
\omterm{def:b2:proba:space}{kejadian} kepunahan yang \omterm{def:b2:proba:independence}{saling bebas}: peluangnya $q^k$. Untuk
$\lambda = 1.5$: peluang wabah $1 - q^k \geq 0.99$
menuntut $q^k \leq 0.01$, yakni $k \geq
\frac{\ln 0.01}{\ln 0.417} \approx 5.3$: jadi enam kasus awal
membuat wabahnya $99\%$ pasti.

\textbf{24.} \emph{$G'(q) < 1$:} $G - \mathrm{id}$ \omterm{def:b2:affine:convex}{cembung}
dan lenyap di $q$ dan $1$, jadi ia $\leq 0$ pada $\intcc
q1$; bila $G'(q) = 1$, maka singgungnya di $q$ (yang oleh kecembungannya
diletakkan di bawah $G$) akan memaksa $G(t) \geq t$ pada $\intcc q1$,
sehingga $G \equiv \mathrm{id}$ di sana, yang mematikan semua koefisien
$p_n$ ($n \geq 2$) dan bertentangan dengan $m > 1$. \emph{Kekonvergenan
geometriknya:} $q_n < q$ untuk semua $n$ (induksi, dengan $G$
naik), dan teorema nilai rata-rata memberi $q - q_{n+1} =
G'(c_n)(q - q_n)$ dengan $c_n \in \intoo{q_n}q$, jadi $G'(c_n)
\leq G'(q) < 1$ dan $q - q_n \leq q\,G'(q)^n$. \emph{Proses
pendampingnya:} $\widehat G(t) = G(qt)/q =
\sum_kp_kq^{k-1}t^k$ berkoefisien taknegatif dan
$\widehat G(1) = G(q)/q = 1$: sebuah \omterm{ex:b2:powerseries:fibonacci}{fungsi pembangkit}; reratanya adalah $\widehat
G'(1) = G'(q) < 1$: subkritis. Keluarga geometriknya: $G(t) =
\frac{q}{1-pt}$, $q_{\mathrm{ext}} = \frac qp$, dan
\[
\widehat G(t) = \frac pq\cdot\frac{q}{1 - p\frac qp t}
= \frac{p}{1 - qt} :
\]
yakni \omterm{def:b2:randomvar:law}{hukum} keturunan geometrik dengan $p$ dan $q$ tertukar ---
proses superkritis yang dipandang pada \omterm{def:b2:proba:space}{kejadian} kepunahannya adalah
cerminan subkritisnya.

\textbf{25.} Tabelnya: $m < 1$: $q = 1$, $\P(Z_n > 0)
\asymp m^n$ (pertanyaan 4--5), $\E(Y) = \frac1{1-m}$, dan
generasi sintas dengan rerata bersyarat yang terbatas. $m = 1$:
$q = 1$, $\P(Z_n > 0) \sim \frac2{\sigma^2n}$ (Kolmogorov),
$\E(Y) = \infty$ dengan $\P(Y > n) \asymp n^{-1/2}$, dan penyintas
berukuran $\sim \frac{\sigma^2n}2$. $m > 1$: $q < 1$ merupakan
titik tetap terkecilnya, $q - q_n = O(G'(q)^n)$, pertumbuhan $\E(Z_n) = m^n$,
dan dengan disyaratkan pada kematiannya prosesnya adalah
pendamping subkritisnya (pertanyaan 24). Adapun perkakasnya: komposisi
\omterm{ex:b2:powerseries:fibonacci}{fungsi pembangkit} mengubah rekursi populasi menjadi iterasi fungsi;
kecembungannya memancangkan geometri titik tetapnya; Taylor di
$1^-$ mengubah hipotesis momen menjadi uraian lokal; dan
perata-rataan Cesàro memeras $1/n$ Kolmogorov dari sebuah
jumlah teleskopik. Adapun jilid Tahun ke-3 menambahkan martingal $Z_n/
m^n$ --- yang limit hampir pastinya menghaluskan $\E(Z_n) = m^n$ menjadi
laju tumbuh lintasan demi lintasan --- dan teorema
Yaglom, yakni \omterm{def:b2:randomvar:law}{hukum} limit di balik geometri bersyarat yang
teramati pada pertanyaan 11.

\end{solution}
